\renewcommand{\algorithmcfname}{方法}

%%%%%%%%%%%%%%%%%%%%%%%%%%%%%%%%%%%%%%%%%%%%%%%%%%%%%%%%%%%%%%%%%%%%%%%
\chapter{一种基于多策略检索路由的 RAG 方法与实验}\label{chap03}

\section{引言}\label{chap03:intro}

\subsection{问题定义与分析}\label{chap03:problem_def}

检索增强生成（RAG）系统的核心挑战，本质上在于如何在生成质量与计算开销之间建立最佳的平衡机制。从系统论的角度，我们将这一自适应路由问题形式化为一个典型的约束优化问题。

假设输入查询空间为 $\mathcal{Q}$，系统的检索策略空间为 $\mathcal{S}$。策略空间 $\mathcal{S}$ 包含了一系列离散的动作，如直接利用参数化知识生成（Zero-Retrieval）、单次检索（Single-Hop）、以及迭代式多跳检索（Multi-Hop）等。对于任意输入查询 $q \in \mathcal{Q}$，系统的目标是寻找最优策略 $s^* \in \mathcal{S}$，使得最终的生成质量 $Q(y|q, s)$ 最大化，同时将系统整体开销 $C(q, s)$（包括时间延迟、计算资源消耗及潜在的噪声干扰）控制在最优范围内。这构成了帕累托最优（Pareto Optimality）的求解过程：
\begin{equation}
s^* = \mathop{\arg\max}_{s \in \mathcal{S}} \left[ \lambda \cdot Q(y|q, s) - (1-\lambda) \cdot C(q, s) \right]
\end{equation}
其中，$\lambda$ 为平衡生成质量与系统开销的权衡系数。

为了求解上述方程，必须首先认识到不同查询 $q$ 对策略 $s$ 的需求存在本质差异。如表 \ref{tab:query_complexity} 所示，查询的逻辑复杂度直接决定了其所需的证据深度。对于低复杂度问题（$L_0$），强制检索不仅增加开销 $C$，更可能引入噪声从而降低 $Q$；而对于高复杂度问题（$L_2$），浅层检索则会导致证据链缺失。因此，传统的静态 RAG 系统对所有 $q$ 采用固定 $s$ 的做法，显然无法在动态变化的查询流中实现全局最优。

\begin{table}[htbp]
    \centering
    \caption{不同查询复杂度下的检索需求特征与策略映射}
    \label{tab:query_complexity}
    \renewcommand{\arraystretch}{1.3}
    \begin{tabular}{l p{3.5cm} l p{4.5cm}}
        \toprule
        \textbf{复杂度} & \textbf{典型示例} & \textbf{最优策略 ($s^*$)} & \textbf{非最优策略的代价 ($C$)} \\
        \midrule
        \textbf{$L_0$} & 逻辑推理、代码生成、\newline 闲聊对话 & 参数化生成 ($Z$) & \textbf{过度检索}：引入语义噪声 ($\mathcal{N}$)，导致推理退化与资源浪费 \\
        \textbf{$L_1$} & 单一实体属性查询、\newline 简单事实问答 & 混合检索 ($S$) & \textbf{精度受限}：单一召回模式可能无法覆盖语义或词法匹配需求 \\
        \textbf{$L_2$} & 多实体关系比较、\newline 跨文档推理 & 逻辑引导推理 ($M$) & \textbf{检索不足}：缺乏完整证据链导致推理中断，引发事实性幻觉 \\
        \bottomrule
    \end{tabular}
\end{table}

为了逼近公式 (3.1) 中的最优解，近期的研究开始探索自适应检索机制 \cite{asai2024selfrag, jeong2024adaptive}。然而，现有的自适应方法主要采取监督学习范式，试图训练一个判别模型来拟合生成模型的“知识边界”，即预测“LLM 此时此刻是否知道答案”。本章认为，这种基于“隐式知识边界预测”的求解路径存在严重的理论缺陷，导致其难以鲁棒地实现目标：

首先，存在显著的模型容量失配问题。作为路由器的判别模型通常是轻量级模型，如 BERT-base模型，而生成模型往往是拥有千亿参数的超大模型。试图用一个小容量模型去“记住”或“覆盖”大容量模型所蕴含的庞大参数化知识边界，在信息论上是一个高损压缩过程。小模型无法精准地模拟大模型的认知盲区，极易导致对 $Q(y|q, s)$ 的错误预估。

其次，面临静态与动态知识解耦（Static-Dynamic Decoupling）的挑战。大语言模型的知识通过持续预训练或微调不断更新，而作为控制器的判别模型一旦训练完成，其对“知识边界”的认知即被冻结。这种非同步性会导致严重的认知错位：当 LLM 已经掌握了新的事实知识时，判别模型仍可能基于旧的训练数据判定其“未知”并强制检索，从而引入不必要的噪声，导致系统开销 $C(q,s)$ 不降反升。

基于上述分析，本章提出一种新的求解思路：高效的检索路由不应依赖于对 LLM “知不知道”的隐式预测，而应依赖于对问题本身“难不难”的显式逻辑判断。问题的逻辑结构是客观存在的固有属性，不随 LLM 的版本迭代而改变。因此，本章旨在构建一个基于问题复杂度感知的自适应路由框架，通过将连续的复杂度特征映射至离散的最优动作空间，以打破效率与准确率之间的刚性权衡。

\subsection{方法总览与章节结构}\label{chap03:chapter_structure}

\begin{figure}[htbp]
    \centering
    \includegraphics[width=15cm]{imgs/chapter03/overview.pdf}
    \caption{基于复杂度感知与逻辑引导的自适应 RAG 路由框架整体结构图}
    \label{fig:framework}
\end{figure}

为系统性地解决检索增强生成中“效率—准确率”难以兼顾的问题，本章提出了一种基于复杂度感知与逻辑引导的自适应 RAG 路由框架。如图 \ref{fig:framework} 所示，该框架整体采用“感知—决策—执行”（Perception--Decision--Execution, PDE）的分层控制范式，将传统 RAG 系统中高度耦合的检索与生成流程解构为若干功能清晰、职责明确的模块，从而实现对不同复杂度查询的精细化调度。

图 \ref{fig:framework} 展示了系统的整体数据流与控制流结构。系统主要由三类核心组件构成：检索基础设施层、复杂度感知与路由决策层，以及差异化执行层。用户查询首先进入复杂度感知模块，经由语义路由分类器完成对问题结构与检索需求的联合判别，随后由动作分发器选择最优执行路径，最终在对应的检索与推理配置下完成答案生成。

在系统底层，检索基础设施层为上层模块提供统一、可组合的检索能力支持。如图左侧所示，系统同时构建了稀疏检索与稠密检索两类索引结构，并在此基础上支持混合检索与重排序流程。其中，稀疏检索以 BM25 与 SPLADE 为代表，侧重于词法匹配与高精度召回；稠密检索采用 bi-encoder架构 与 HNSW 索引结构，能够更好地覆盖语义相似但表述差异较大的查询场景。该层的设计为不同复杂度问题提供了稳定的能力原语。

在检索基础设施之上，系统引入复杂度感知与路由决策机制，如图中上部所示。与现有依赖“知识边界预测”的方法不同，本章从问题本身出发，对查询所隐含的推理复杂度与检索形态需求进行显式建模。路由分类器以用户查询为输入，输出位于动作空间
\[
\mathcal{A} = \{Z, S_{\text{sparse}}, S_{\text{dense}}, S_{\text{hybrid}}, M\}
\]
中的离散决策结果，并作为系统的全局控制信号。分类器的构建、特征建模方式及复杂度分层定义，将在第 \ref{chap03:perception} 节中进行系统阐述。

完成动作决策后，系统通过动作分发器将查询路由至不同的执行分支，对应图 \ref{fig:framework} 下部所示的多条差异化执行路径。对于被判定为低复杂度的查询，系统采用零检索路径，直接旁路检索模块，仅依赖语言模型的参数化知识进行推理与生成，以最小化系统延迟与噪声引入。对于中等复杂度问题，系统触发单跳检索路径，根据路由结果动态选择稀疏、密集或混合检索方式，并在一次检索后完成上下文注入与答案生成。对于高复杂度查询，系统进一步启用逻辑引导的多跳推理路径，通过显式的逻辑规划模块，将原始问题拆解为若干子目标，并在“规划—检索—推理—校验”的闭环过程中逐步构建完整证据链。

上述不同执行路径在检索次数、上下文构建方式以及推理控制粒度上均存在显著差异，其设计动机与具体实现细节将在第 \ref{sec:execution} 节中展开讨论。

基于上述整体框架设计，本章后续内容将围绕各功能模块依次展开。第 \ref{chap03:perception} 节重点讨论问题复杂度建模与路由分类体系；第 \ref{sec:execution} 节介绍自适应 RAG 的差异化执行策略，包括零检索、单跳混合检索以及逻辑引导多跳推理机制；随后在第 \ref{sec:experiment} 节与第 \ref{sec:results} 节中，从分类性能与端到端问答效果两个层面，对所提方法进行系统实验验证；最后在第 \ref{chap03:summary} 节对本章方法进行总结，并讨论其在复杂问题问答场景中的适用性与扩展潜力。

通过上述章节组织结构，本章旨在从问题建模、路由决策与执行控制三个层面，系统论证复杂度感知自适应 RAG 框架在理论合理性与实践有效性上的优势。

\section{问题复杂度感知与路由决策建模}\label{chap03:perception}

\subsection{自适应路由问题的形式化建模}

为了在保证生成质量的前提下实现检索资源的按需分配，我们将 rag 系统中的路由决策过程形式化为一个受约束的分层分类问题。设查询空间为 $\mathcal{q}$，系统可执行的动作空间为 $\mathcal{a}$。路由器的目标是学习一个映射函数
\[
f_\theta: \mathcal{q} \rightarrow \mathcal{a},
\]
使得对于任意查询 $q \in \mathcal{q}$，所选动作 $a \in \mathcal{a}$ 能够在最大化生成效果的同时控制整体计算开销。

不同于端到端的黑盒映射方式，我们引入一组显式的中间语义变量，将路由决策拆解为多个受约束的子判别问题，从而提升模型决策的可解释性与鲁棒性。具体而言，我们构建了一个由问题复杂度、检索形态与执行动作共同组成的分层路由建模框架。

首先，引入问题复杂度变量 $\mathcal{c}$，用于刻画查询在信息依赖与推理结构上的难度层级：
\[
\mathcal{c} \in \{l_0, l_1, l_2\},
\]
其中，$l_0$ 表示无需外部检索即可完成回答的问题，$l_1$ 表示仅需单跳检索即可获取关键信息的问题，而 $l_2$ 表示需要多跳检索与多步推理的问题。

其次，对于需要检索的查询，引入检索形态变量 $\mathcal{r}$，用于描述检索过程中应侧重的匹配信号类型，包括基于词法匹配的稀疏检索、基于语义相似度的稠密检索，以及二者的混合形式。

在上述中间变量的约束下，系统的执行动作空间 $\mathcal{a}$ 被限制为一组结构化的候选策略，包括不执行检索的直接生成策略 $z$，不同检索形态下的单步检索策略 $s_{\text{sparse}}$、$s_{\text{dense}}$、$s_{\text{hybrid}}$，以及面向多跳问题的迭代式检索与推理策略 $m$。表~\ref{tab:routing_taxonomy} 总结了问题复杂度、检索形态与执行动作之间的对应关系。

通过上述分层建模，原本复杂的查询到动作映射被分解为若干受语义约束的子决策，为后续的表示学习分析与系统实例化提供了清晰的结构基础。

\subsection{基于实体结构的隐式表示学习机制}

在上述分层建模框架中，问题复杂度 $\mathcal{c}$ 与检索形态 $\mathcal{r}$ 并非通过显式规则或人工特征进行判定，而是由模型从查询文本中隐式预测。本节从表示学习的角度，分析模型如何基于查询中的实体结构与语义模式，自动习得对路由相关中间变量的判别能力。

一个关键观察是，查询中所包含的实体分布特征与实体间的关系结构，能够有效反映问题在信息依赖层面的复杂程度。对于被判定为 $l_1$ 的单跳检索问题，其语义结构通常可抽象为原子化的“实体–属性”对，即查询聚焦于单一核心实体 $e$ 的某一明确属性 $p$。这类查询在结构上呈现出低实体密度与浅关系深度的特征。模型通过注意力机制与上下文建模，能够识别这种近似星型的局部结构，从而预测其仅需一次外部检索即可完成回答。

相比之下，$l_2$ 类多跳问题通常涉及多个实体集合 $\{e_1, e_2, \dots, e_n\}$，且实体之间存在隐式的桥接、比较或嵌套关系。例如，“$e_1$ 出生地的首都是什么？”这一查询隐含了 $e_1 \rightarrow \text{location} \rightarrow \text{capital}$ 的推理链。模型在训练过程中逐渐学会通过上下文中的关系提示词与语义依赖模式，感知查询的关系深度（relation depth），从而将其映射至需要多步信息聚合的复杂度层级。

在完成对问题复杂度的隐式判别后，模型进一步针对 $l_1$ 类查询学习检索形态 $\mathcal{r}$ 的选择逻辑。该过程同样并非依赖显式规则，而是内化为模型参数中的隐式映射关系。当查询包含高区分度的专有名词、代码标识符或精确术语时，词法匹配信号往往更具判别力，模型因此倾向于激活基于倒排索引的稀疏检索策略；当查询以自然语言描述或概念性表述为主、缺乏显著关键词时，模型则更依赖语义嵌入的相似性，选择稠密向量检索以提升召回覆盖；而在多实体共现或约束条件不完整的情况下，模型会预测混合检索策略能够通过稀疏与稠密结果的互补性，在精确匹配与语义泛化之间取得更优平衡。

总体而言，模型通过端到端的监督学习，在隐层空间中形成了对实体分布特征、关系结构复杂度以及语义粒度的连续表征。这一表示为后续将复杂度预测与检索策略选择组织为统一的路由决策流程提供了必要的语义基础。

\subsection{两阶段路由决策流程与策略空间}

在前两节中，我们分别从形式化建模与表示学习的角度定义并分析了路由决策所涉及的中间变量。本节将问题复杂度 $\mathcal{c}$ 与检索形态 $\mathcal{r}$ 显式组织为一个可执行的两阶段路由决策流程:

整体路由过程可划分为两个顺序执行的阶段。第一阶段以查询 $q$ 为输入，对其问题复杂度进行判别，即预测 $\mathcal{c} \in \{l_0, l_1, l_2\}$。当查询被判定为 $l_0$ 时，系统直接采用无需外部检索的生成策略 $z$；当查询被判定为 $l_2$ 时，则触发面向多跳推理的迭代式检索策略 $m$。仅当查询被判定为 $l_1$ 时，系统才进入第二阶段。

第二阶段针对 $l_1$ 类查询，进一步预测其最优检索形态 $\mathcal{r}$，并据此选择对应的单步检索策略，包括基于稀疏匹配的 $s_{\text{sparse}}$、基于语义相似度的 $s_{\text{dense}}$，以及融合二者的 $s_{\text{hybrid}}$。这种设计使得系统在保证决策灵活性的同时，避免了不必要的检索与推理开销。

\begin{table}[htbp]
    \centering
    \small
    \caption{问题复杂度、检索形态与执行动作的策略空间映射}
    \label{tab:strategy_mapping}
    \renewcommand{\arraystretch}{1.3}
    \begin{tabular}{c c c}
        \toprule
        \textbf{问题复杂度 $\mathcal{c}$} &
        \textbf{检索形态 $\mathcal{r}$} &
        \textbf{执行动作 $\mathcal{a}$} \\
        \midrule
        $l_0$ & none   & $z$ \\
        
        $l_1$ & sparse & $s_{\text{sparse}}$ \\
        $l_1$ & dense  & $s_{\text{dense}}$ \\
        $l_1$ & hybrid & $s_{\text{hybrid}}$ \\
        
        $l_2$ & hybrid (iterative) & $m$ \\
        \bottomrule
    \end{tabular}
\end{table}


\begin{table}[htbp]
    \centering
    \small
    \caption{典型查询示例及其对应的路由决策路径}
    \label{tab:routing_examples}
    \renewcommand{\arraystretch}{1.35}
    \begin{tabular}{p{7.2cm} c c c p{4.2cm}}
        \toprule
        \textbf{查询示例} &
        $\mathcal{c}$ &
        $\mathcal{r}$ &
        $\mathcal{a}$ &
        \textbf{路由路径} \\
        \midrule
        
        电影《the decision of christopher blake》的导演去世于哪里？ 
        & $l_1$ & sparse & $s_{\text{sparse}}$
        & $q \rightarrow l_1 \rightarrow \text{sparse}$ \\
        
        哈维尔·埃尔南德斯（小豌豆）与安东尼奥·瓦伦西亚和迪米塔尔·贝尔巴托夫一同进球的那场比赛在哪里举行？
        & $l_2$ & hybrid & $m$
        & $q \rightarrow l_2 \rightarrow m$ \\
        
        与那部关于彼得·特里富诺维奇去世城市的电影具有相同共同官方语言的国家，是什么时候派出独立代表队参加奥运会的？
        & $l_2$ & hybrid & $m$
        & $q \rightarrow l_2 \rightarrow m$ \\
        
        是谁提出了自然选择过程？
        & $l_0$ & none & $z$
        & $q \rightarrow l_0 \rightarrow z$ \\
        
        皇后乐队是哪一年入选英国名人堂的？
        & $l_1$ & sparse & $s_{\text{sparse}}$
        & $q \rightarrow l_1 \rightarrow \text{sparse}$ \\
        
        按陆地面积计算，最大的西班牙语国家是哪一个？
        & $l_1$ & dense & $s_{\text{dense}}$
        & $q \rightarrow l_1 \rightarrow \text{dense}$ \\
        
        \bottomrule
    \end{tabular}
\end{table}

表\ref{tab:strategy_mapping} 系统性地刻画了问题复杂度 $\mathcal{c}$、检索形态 $\mathcal{r}$ 与最终执行动作 $\mathcal{a}$ 之间的映射关系。该映射并非穷举式的策略组合，而是在前述分层建模约束下形成的结构化策略空间：问题复杂度作为第一层决策变量，对可选动作集合施加了强约束，而检索形态仅在单跳检索场景（$l_1$）下被进一步细分。通过这种方式，系统在保证策略覆盖性的同时，有效避免了在低复杂度问题上触发不必要的检索或推理流程，从而实现生成质量与计算效率之间的平衡。

进一步地，表\ref{tab:routing_examples} 通过来自真实数据集的典型查询示例，具体展示了所提出路由机制在实际应用中的决策行为。可以看到，不同查询在实体数量、关系结构以及语义表达方式上的差异，会导致模型在复杂度判定阶段做出不同的预测结果，并由此触发差异化的执行路径。对于无需外部知识支持的概念性问题，系统能够直接路由至无需检索的生成策略；而对于涉及多实体关联或隐式推理链的问题，则会被一致地判定为高复杂度查询，并交由迭代式检索与推理策略处理。

\begin{algorithm}[htbp]
\caption{复杂度感知的自适应 RAG 总体流程}
\label{alg:adaptive_rag}
\KwIn{用户查询 $q$}
\KwOut{生成答案 $y$}

\textbf{/* Perception Stage */} \\
$\mathcal{c} \gets$ PredictComplexity$(q)$ \;

\eIf{$\mathcal{c} = l_0$}{
    \textbf{/* Zero-Retrieval Execution */} \\
    $y \gets \mathcal{G}(q)$ \;
}{
    \eIf{$\mathcal{c} = l_1$}{
        \textbf{/* Single-Hop Decision */} \\
        $\mathcal{r} \gets$ PredictRetrievalType$(q)$ \;
        $D \gets$ Retrieve$(q, \mathcal{r})$ \;
        $C \gets$ RerankAndSelect$(q, D)$ \;
        $y \gets \mathcal{G}(q, C)$ \;
    }{
        \textbf{/* Multi-Hop Reasoning (M-Action) */} \\
        Initialize reasoning state $\mathcal{S}$ \;
        $y \gets$ MultiHopReasoning$(q, \mathcal{S})$ \;
    }
}
\Return $y$
\end{algorithm}

 伪代码\ref{alg:adaptive_rag} 对前述复杂度感知与两阶段路由决策机制进行了统一的过程化刻画，展示了系统如何以查询 $q$ 为起点，依次完成问题复杂度判定、检索形态选择以及执行策略调度。通过将路由决策显式组织为“感知—决策—执行”的分层流程，该方法避免了端到端黑盒路由在复杂场景下可能出现的不可控行为，使不同执行策略之间的调用关系具备清晰的语义边界与可解释性。

与现有将检索与生成紧耦合的 RAG 方法不同，本文所提出的路由机制以问题复杂度为核心控制变量，对可选执行动作空间施加结构化约束，从而在保证生成质量的同时，有效抑制了不必要的检索与推理开销。该设计在思想上与近年来提出的基于推理过程调度的检索增强方法一脉相承，但进一步强调了复杂度感知在整体系统控制中的中心作用，为自适应 RAG 系统提供了一种更具可扩展性的决策范式。

在此统一路由框架下，不同执行策略（零检索、单跳检索与多跳推理）被视为可插拔的功能模块，其内部实现彼此解耦。下一节将分别对混合检索策略的理论设计，以及单跳与多跳执行策略的具体实现机制进行详细说明，以完整展示所提出自适应 RAG 框架在不同复杂度查询下的执行行为。

\section{自适应 RAG 执行策略设计}

\subsection{混合检索策略理论设计}

单一的检索模态在处理异构查询时往往存在固有的局限性：稀疏检索（Sparse Retrieval）依赖于词汇重叠，难以捕捉语义同义性（Vocabulary Mismatch）；而密集检索（Dense Retrieval）虽然擅长语义匹配，却容易忽略精确实体信息或产生语义漂移。为了兼顾召回率（Recall）与精确率（Precision），本系统设计了一种基于正交特征互补的混合检索策略，作为单跳与多跳策略底层的核心支撑。

混合检索的理论框架建立在多阶段排序（Multi-Stage Ranking）范式之上，其执行流程可形式化为“候选生成-特征融合-精细重排”的级联结构：

\begin{enumerate}
    \item \textbf{候选生成阶段}：系统并行执行基于 BM25 的词法检索与基于 HNSW 的向量检索，分别获取候选文档集合 $\mathcal{D}_{lex}$ 与 $\mathcal{D}_{sem}$。
    \item \textbf{特征融合阶段}：为解决不同检索器打分分布的异质性问题，引入基于分布归一化的线性加权融合机制。对于任意文档 $d$ 与查询 $q$，其融合得分 $S_{\text{fusion}}$ 定义为：
    \begin{equation}
    S_{\text{fusion}}(q, d) = \alpha \cdot \mathcal{N}(S_{\text{lexical}}(q, d)) + (1-\alpha) \cdot \mathcal{N}(S_{\text{semantic}}(q, d))
    \end{equation}
    其中，$\mathcal{N}(\cdot)$ 为 Min-Max 归一化函数，$\alpha$ 为调节权重的超参数，可视作问题复杂度 $C(q)$ 的函数 $\alpha(C(q))$。
    \item \textbf{精细重排阶段}：对融合后的 Top-$K$ 候选集引入 Cross-Encoder 模型进行深度全注意力交互，以捕捉细微的语义关联并滤除噪声。
\end{enumerate}

该混合策略在理论上构建了鲁棒的检索边界：利用稀疏检索保证关键实体的精确命中，利用密集检索提升语义覆盖，最后通过重排序优化最终候选集质量。该理论是后续单跳与多跳策略实现高质量上下文支持的基础。

\subsection{零检索与单跳策略}

针对低复杂度的查询请求，系统设计了轻量化的执行路径，以最小计算开销获取准确响应。根据前序分类器的决策，具体分为零检索与单跳检索两种基本模式。

\textbf{零检索策略（Z）} 主要针对 $L_0$ 级别的查询（如逻辑谜题、代码生成或日常闲聊）。系统完全旁路检索模块，直接将原始查询 $q$ 输入生成器 $\mathcal{G}$。该策略基于两点假设：一是大语言模型的参数化知识 $\theta$ 足以处理此类通用任务；二是非知识密集型任务中，强制检索可能引入噪声，导致注意力分散或幻觉产生。Z 策略显著降低系统延迟，并避免外部噪声干扰。

\textbf{单跳检索策略（S）} 针对 $L_1$ 级查询，需要外部事实支撑。系统构建统一检索后端，集成多种异构算法：基于 HNSW 的 Dense Retrieval 捕捉语义相似性；基于 SPLADE 的学习型稀疏检索弥合词法与语义鸿沟。单跳策略中引入“黄金三角”检索-重排机制：
\begin{enumerate}
    \item \textbf{召回阶段}：利用大模型提取关键实体生成短查询 $q_{\text{short}}$，最大化召回范围，确保候选文档集 $\mathcal{D}_{\text{cand}}$ 包含潜在相关信息。
    \item \textbf{重排阶段}：使用原始长查询 $q_{\text{long}}$ 配合 Cross-Encoder 对 $\mathcal{D}_{\text{cand}}$ 进行精细重排序，实现“短查询广召回、长查询精重排”。
\end{enumerate}

\subsection{逻辑引导的多跳推理策略}

针对多跳问答任务中复杂知识依赖、跨文档证据整合以及推理路径不可控等问题，本文提出一种逻辑引导的多跳推理策略（Multi-hop Action Strategy, M-Action）。该策略不再将多跳推理视为固定长度的线性推理链，而是将其形式化为一个由大语言模型驱动、以外部知识为约束的状态化决策过程，使模型能够在结构化约束下进行逐步检索与推理。整体框架如图 \ref{fig:m_action_framework} 所示。

在系统初始化阶段，首先对原始问题 $q$ 进行逻辑规划。具体而言，大语言模型根据问题语义生成高层逻辑计划 $\pi$，用于刻画问题的潜在推理路径及子目标之间的逻辑依赖关系。逻辑计划不仅为后续多跳检索提供方向性约束，同时作为全局推理策略嵌入到后续推理提示中，从而引导模型避免无目标搜索。随后，系统基于原始问题 $q$ 执行初始检索，以获得与问题直接相关的候选证据，从而实现推理过程的“热启动”。

在多跳推理过程中，系统在每一轮迭代中维护一个全局推理状态，用于刻画当前推理进展与知识积累情况。形式化地，在第 $t$ 轮推理时，定义推理状态为
\begin{equation}
\mathcal{S}_t = \left(\pi, H_t, \mathcal{E}_t\right),
\end{equation}
其中，$\pi$ 表示逻辑计划，$H_t$ 表示截至当前轮次的推理历史，记录模型已执行的检索查询与决策轨迹，$\mathcal{E}_t$ 表示当前已获取的证据集合，由外部知识库检索得到并经过筛选的文本片段构成。推理状态的显式建模使得多跳推理过程具备可追踪性与可控性。

在状态 $\mathcal{S}_t$ 下，模型需要选择下一步动作。本文将模型的决策空间限定为有限动作集合
\begin{equation}
a_t \in \mathcal{A} = \{\textsc{Retrieve}, \textsc{Answer}\},
\end{equation}
其中，$\textsc{Retrieve}$ 表示发起新的检索请求，$\textsc{Answer}$ 表示基于当前证据集合生成候选答案。动作选择由大语言模型完成，其决策依据主要来源于当前推理状态 $\mathcal{S}_t$ 与逻辑计划 $\pi$。当模型判断当前证据集合尚不足以支持问题求解时，倾向于选择检索动作；当模型认为已有证据能够支持答案生成时，则选择回答动作。

当执行检索动作时，模型根据当前推理状态生成检索查询 $q_t$。为降低模型生成查询时的噪声与冗余，系统首先对查询进行规范化处理，包括对话式冗余短语消除、实体保留与语义去重等操作，从而得到语义更集中的检索表达。随后，系统通过混合检索策略从外部知识库中获取候选文档集合：
\begin{equation}
\mathcal{D}_t = \mathrm{Retrieve}(q_t),
\end{equation}
其中检索模块融合稀疏检索与稠密检索结果，并通过重排序模型进行筛选。为避免多跳推理过程中出现重复检索或循环检索，系统进一步引入语义相似度检测机制，对当前查询与历史查询进行比对，从而阻止语义高度相似的查询再次执行。

在获得候选文档后，系统对检索结果进行过滤与去重，将有效证据写入证据集合，并更新为
\begin{equation}
\mathcal{E}_t \leftarrow \mathcal{E}_t \cup \mathrm{Filter}(\mathcal{D}_t).
\end{equation}
与此同时，推理历史被更新为
\begin{equation}
H_{t+1} = H_t \cup \{(q_t, \mathcal{D}_t)\},
\end{equation}
从而显式刻画知识在多跳推理过程中的累积式演化。

当模型选择回答动作时，系统基于当前证据集合 $\mathcal{E}_t$ 生成候选答案 $\hat{y}$。为了避免模型在证据不足或信息缺失的情况下产生幻觉，系统引入证据一致性校验机制，用于约束回答动作的触发条件。形式化地，定义证据支持函数为
\begin{equation}
\mathcal{V}(\hat{y}, \mathcal{E}_t) \in \{0,1\},
\end{equation}
其中，当 $\mathcal{V}(\hat{y}, \mathcal{E}_t)=1$ 时，表示当前答案与证据集合之间不存在明显冲突，且未出现“未找到相关信息”等否定信号，系统允许进入最终答案生成阶段；否则，系统拒绝当前回答，并重新触发检索动作以获取更多相关证据。需要指出的是，该证据校验机制并非严格的逻辑证明，而是一种宽松的语义一致性约束，其实现方式与实际系统中的证据过滤策略保持一致。

在实际推理过程中，系统还引入多重控制机制以增强推理稳定性。首先，通过对语言模型输出进行物理截断，避免模型在单次推理中生成过长的多跳推理链，从而减少幻觉传播风险。其次，在每一轮推理中，系统允许模型在有限次数内重新生成动作决策，以提升决策的鲁棒性。当模型连续多次生成无效查询或出现格式错误时，推理过程将被提前终止。此外，系统对推理步数设置上限，以避免推理过程无限循环。

综合上述机制，多跳推理过程可以被视为一个由状态驱动的闭环决策系统，其状态转移过程可表示为
\begin{equation}
\mathcal{S}_t \xrightarrow{a_t} \mathcal{S}_{t+1},
\end{equation}
模型通过不断更新推理状态与证据集合，逐步逼近问题的最终答案。当模型选择回答动作且通过证据校验，或达到最大推理步数限制时，推理过程终止并输出最终答案。

与传统 RAG 框架相比，M-Action 不仅引入逻辑规划与状态建模机制，而且通过显式动作空间与证据约束，使语言模型的推理轨迹受到结构化控制，从而在复杂多跳推理场景下显著提升系统的可控性、稳定性与可解释性，并有效缓解无约束生成所导致的事实性幻觉问题。

\section{实验设计}\label{ch03:exp}

\subsection{数据集构建与预处理}
\label{sec:dataset_construction}

\subsubsection{端到端评估数据集选取}

为了系统评估自适应路由框架在不同问题复杂度场景下的性能表现，本文选取了六个具有代表性的开放域问答数据集作为基准测试集。这些数据集覆盖了从单一事实检索到多跳推理的不同任务类型，能够较为全面地反映端到端 RAG 系统在不同推理深度下的行为特征。表 \ref{tab:datasets} 总结了各数据集的基本特征与侧重点。

根据问题所需的推理深度，我们将评估数据集划分为单跳任务与多跳任务两类。在单跳事实类任务中，选用了 SQuAD v2.0、Natural Questions (NQ) 与 TriviaQA。SQuAD v2.0 主要考察模型对文本片段的精确抽取能力；NQ 来源于真实用户搜索查询，具有较强的语言多样性与开放域特征；TriviaQA 则包含更复杂的实体关系与长文本语境。这类数据集主要用于评估系统在单次检索场景下的检索与生成能力，同时检验路由模型是否能够避免对简单问题触发不必要的复杂推理策略。

在多跳推理类任务中，本文采用了 HotpotQA、2WikiMultiHopQA 与 Musique。HotpotQA 要求模型通过桥接实体完成跨文档推理；2WikiMultiHopQA 涉及多实体的属性比较与信息聚合；Musique 则强调多步组合推理，其问题设计显著降低了单跳检索的可行性。这些数据集构成了对系统多跳推理策略的主要评估场景，用于分析模型在复杂知识依赖条件下的规划与执行能力。

为了保证不同数据集之间的可比性，本文对评估数据进行了统一的标准化处理。具体而言，我们从各数据集的官方开发集中进行分层抽样，构建样本规模一致的评估子集，并对问题文本、答案格式与检索语料结构进行统一规范。每条样本包含唯一标识符、问题文本、标准答案对象以及包含支持性事实的上下文列表，以消除不同数据源格式异构带来的评估障碍。该标准化过程旨在消除不同数据集在规模、文本风格与标注形式上的差异，使其能够在同一端到端评估框架下进行横向对比，从而更准确地刻画路由机制在不同复杂度场景下的性能表现。


\begin{enumerate}
\item \textbf{多跳推理数据集（Multi-hop Datasets）}：此类数据集旨在评估系统处理复杂查询的能力，对应本研究定义的 L2 复杂度（需多跳检索）。
\begin{itemize}
\item \textbf{HotpotQA (Bridge Type)}：该数据集要求模型在多个文档间建立桥梁实体以推导答案。本研究主要选取其 Bridge 类型样本，以测试系统在非显式关联文档间的推理能力。
\item \textbf{Musique (Compositional)}：这是一个具有高度组合性的多跳数据集，其问题通常需要分解为多个逻辑步骤。选取该数据集旨在验证多跳 Agent 在长链条推理中的稳定性。
\item \textbf{2WikiMultiHopQA}：该数据集涉及多个实体间的关系推理，且通常包含比较类问题。它用于测试系统在多实体场景下的检索精度与逻辑综合能力。
\end{itemize}
\item \textbf{单跳事实数据集（Single-hop Datasets）}：此类数据集主要测试系统对单一明确事实的检索能力，对应本研究定义的 L1 复杂度（需单跳检索）。
\begin{itemize}
    \item \textbf{Natural Questions (NQ)}：来源于真实的 Google 搜索查询，问题通常简短且目标明确，用于评估系统对开放域事实问题的响应速度与准确率。
    \item \textbf{SQuAD (Stanford Question Answering Dataset)}：经典的阅读理解数据集，问题通常依赖于单一文档段落。本实验使用其开放域版本设置。
    \item \textbf{TriviaQA}：包含由琐事爱好者编写的复杂事实性问题，虽然部分问题具有一定的描述性，但大多数仍可通过单次检索解决。
\end{itemize}

\end{enumerate}



\begin{table}[htbp]
    \centering
    \caption{端到端评估所采用的基准数据集统计与特征}
    \label{tab:datasets}
    \renewcommand{\arraystretch}{1.2}
    \resizebox{\textwidth}{!}{
    \begin{tabular}{l l l l}
        \toprule
        \textbf{数据集} & \textbf{复杂度类别} & \textbf{推理类型} & \textbf{关键特征描述} \\
        \midrule
        \textbf{SQuAD v2.0} & 单跳 ($L_1$) & 阅读理解 & 包含无法回答的问题，考察精确的文本片段抽取 \\
        \textbf{Natural Questions} & 单跳 ($L_1$) & 开放域检索 & 真实用户搜索查询，覆盖广泛的维基百科事实 \\
        \textbf{TriviaQA} & 单跳 ($L_1$) & 事实问答 & 包含复杂的词汇表述与长文档上下文 \\
        \midrule
        \textbf{HotpotQA} & 多跳 ($L_2$) & 桥接推理 & 需寻找中间实体以连接起止点 \\
        \textbf{2WikiMultiHopQA} & 多跳 ($L_2$) & 比较与聚合 & 涉及多实体的属性比较或合成推理 \\
        \textbf{Musique} & 多跳 ($L_2$) & 组合推理 & 需要多步有序检索才能获得完整证据 \\
        \bottomrule
    \end{tabular}
    }
\end{table}


\subsubsection{路由分类器训练数据构建}

与端到端评估数据集不同，路由模型的训练需要显式的复杂度与检索策略标签作为监督信号。然而，现有公开问答数据集通常仅提供问题与答案对，缺乏面向检索路径规划的中间标注信息。为此，本文面向路由分类器的训练需求，构建了一套融合弱标签提示与人机协同生成的数据构建流程，用于自动生成包含问题复杂度 $\mathcal{C}$、检索形态 $\mathcal{R}$ 与执行动作 $\mathcal{A}$ 的细粒度标签。

\paragraph{1. 弱标签提示学习策略}

训练数据主要来源于六个主流问答数据集的开发集子集，包括 Natural Questions、SQuAD、TriviaQA 以及 HotpotQA、2WikiMultiHopQA 和 Musique。本文并未直接采用数据集原有的任务标签作为训练目标，而是将其作为弱先验信息引入到半自动化标注过程中。

具体而言，我们构建了一个基于大语言模型的结构化标注器。在标注过程中，数据集来源信息被作为非约束性的上下文提示注入模型输入，用于提供任务背景，但不作为强制标签约束。标注模型需要根据问题文本的\textbf{语言特征}与\textbf{逻辑结构}，综合判断其推理复杂度与检索需求。例如，对于呈现“实体—属性”结构的查询，模型倾向于判定其为单跳检索问题，并进一步区分稀疏检索、稠密检索或混合检索策略；对于涉及多实体关联、比较或隐含推理链的问题，则判定其为多跳推理场景。

通过上述过程，模型为每个问题生成包含复杂度层级（$L_0, L_1, L_2$）、索引策略以及系统执行动作（$Z, S_{\text{sparse}}, S_{\text{dense}}, S_{\text{hybrid}}, M$）的结构化标签。这些标签作为监督信号，用于训练路由分类器以学习问题复杂度感知与检索策略选择能力。相较于直接使用数据集标签，该方法能够在一定程度上缓解原始数据集与实际问题复杂度之间的偏差，从而提高标注结果的适用性。

\begin{itemize}
\item \textbf{复杂度标签 (Complexity)}：L0（无需检索）、L1（单跳检索）、L2（多跳检索）。
\item \textbf{索引策略 (Index Strategy)}：None、Lexical（词法）、Semantic（语义）、Hybrid（混合）。
\item \textbf{执行动作 (Action)}：Z（零检索）、S-Sparse、S-Dense、S-Hybrid、M（多跳 Agent）。
\end{itemize}

\paragraph{2. L0 类别数据合成和校验}
主流的 QA 数据集主要关注基于知识的问答，极度缺乏“无需检索”的样本。若仅使用上述 QA 数据集进行训练，分类器极易陷入“过度检索”的局部最优，即对所有输入都执行检索操作，这违背了本研究提升系统效率的初衷。因此，本文通过生成式方法构建了 $L_0$ 类别数据，以平衡训练集的类别分布。合成样本覆盖逻辑推理、代码生成、创意写作、日常对话、文本处理与常识问答等多种场景，用以模拟真实交互中无需外部知识检索即可完成的任务。

为保证合成数据的质量，本文引入了人机协同的校验机制。所有生成样本首先经过自动化一致性筛查，随后进行人工抽样复核，重点剔除包含显式事实依赖、时效性实体或语义歧义的边界样本。最终，经过清洗的 $L_0$ 样本与 $L_1/L_2$ 标注数据进行合并与随机混洗，构建了类别分布相对均衡、语义覆盖较为全面的路由模型训练数据集。

\paragraph{3. 数据质量校验与清洗}

尽管使用了先进的 LLM 进行标注与合成，数据中仍可能存在错误样本。为了确保训练数据的质量，本研究引入了人工校验机制。我们从标注的 QA 数据中随机抽取 20\% 的样本进行了人工校验。

审查标准包括：复杂度标签是否与问题实际逻辑相符；合成数据是否存在幻觉或格式错误；弱标签提示是否有效地引导了模型，而非导致盲从。对于校验中发现的错误模式，人工矫正相关样本，并重新校验直至人工抽检全部合格。

从每个公开数据集随机采样 500 个样本并经过上述构建流程打乱混合后，最终得到了一个总规模为 3,600 条样本的混合训练数据集。该数据集在复杂度维度上呈现出均衡的分布，能够支持分类器学习到鲁棒的决策边界。为了进行模型训练与超参数调优，我们将该数据集按照 9:1 的比例随机划分为训练集和验证集，其中训练集用于LoRA微调、验证集用于监控训练过程中的 Loss 变化及防止过拟合。


\subsection{基线模型与对比方法}
\label{sec:baselines}

为了验证本文提出的基于分类器的动态路由机制在效率与准确率平衡方面的优势，并系统分析不同检索策略对问答性能的影响，本研究构建了多层次的对比实验体系。对比方法涵盖了从完全依赖模型参数知识的生成方式，到静态检索增强框架，再到现有动态检索方法，从而形成具有代表性的基线体系。

为了保证实验结果的可比性，所有实验设置均采用相同的生成模型与检索后端，仅在检索策略的触发方式与路由机制上存在差异。通过这种控制变量的方式，可以更加准确地评估动态路由策略对系统性能的贡献。

\subsubsection{无检索生成}

在无检索生成设置下，系统不调用任何外部检索模块，而是直接将用户查询输入大语言模型进行回答。该模式对应本文路由策略中的零检索动作，系统完全依赖模型在预训练阶段所内化的知识进行推理与生成。

该设置的主要目的在于建立性能与效率的参考基准。一方面，它刻画了模型在缺乏外部知识支持时的回答能力，从而反映系统在处理长尾事实或专业领域问题时可能出现的幻觉现象；另一方面，由于该模式不引入检索开销，因此能够作为系统响应效率的上限参考。此外，该设置还用于验证逻辑推理、闲聊等无需检索的问题是否能够在闭卷条件下得到有效解决。

\subsubsection{静态检索增强生成}

静态检索增强生成设置模拟了当前工业界较为常见的检索增强范式，即对所有输入问题统一执行检索与生成流程。在该模式下，系统不区分问题类型与复杂度，而是对每个查询均执行混合检索操作，召回一定数量的文档构建上下文，然后由生成模型进行回答。

在检索策略方面，本研究采用词法检索、语义检索与稀疏向量检索的融合方式，并对三种检索结果进行加权整合，从而形成统一的混合检索策略。该设置用于分析两类典型现象：一方面，过度检索可能会向模型引入与问题无关的文档，从而干扰模型对逻辑或常识类问题的判断；另一方面，单次检索在处理复杂推理问题时可能存在信息不足的问题，从而暴露出多跳推理机制的必要性。

\subsubsection{自适应检索增强生成方法复现}

为了与现有动态检索增强生成方法进行对比，本研究复现了一种具有代表性的自适应检索框架作为强基线。该方法通过一个轻量级分类器预测问题复杂度，并据此动态选择检索策略，从而在不同问题类型之间实现策略切换。

在复现过程中，本研究对模型结构与实验条件进行了统一控制，以减少模型能力差异对实验结果的影响。分类器与生成模型均采用与本文方法一致的模型配置，并使用相同的分类提示方式，将问题划分为无需检索、单跳检索与多跳检索三个层级。与本文方法不同的是，该基线方法采用原始数据集提供的标准标签进行训练，并未引入本文提出的弱标签提示与合成数据策略。

该对比设置的核心目的在于剥离模型架构因素，聚焦于训练数据构建策略对路由决策性能的影响，从而评估本文方法在分类器鲁棒性与决策准确性方面的提升效果。

\subsubsection{本文提出的动态路由方法}

本文提出的方法构建了一套基于增强型分类器的动态路由机制。系统首先对输入问题进行复杂度判定，然后在不同检索策略之间进行动态选择。对于无需外部知识即可解决的问题，系统直接进行生成；对于单跳事实性问题，系统执行一次检索并完成回答；对于复杂推理问题，系统调用多跳推理机制，通过多次检索与推理逐步构建答案。

与现有自适应检索方法相比，本文方法的核心差异在于分类器的训练数据构建方式。分类器基于增强数据集进行训练，该数据集不仅包含真实问答样本，还引入了合成的无需检索样本，并通过弱标签提示策略提升标注质量。本研究通过对比实验，验证该数据增强策略在提升分类器对无需检索类别的识别能力方面的作用，并进一步分析其对端到端问答准确率与系统执行效率的综合影响。

\subsection{系统实现与实验配置}
\label{sec:implementation_details}

为了保证实验结果的可复现性与方法评估的客观性，本节从分类器训练策略、生成与检索模块的实现方式以及硬件运行环境三个方面，对系统实现细节与实验配置进行说明。

\subsubsection{分类器训练配置}

路由分类器作为系统决策机制的核心模块，其训练过程需要在模型性能与计算效率之间取得平衡。本研究选取 Qwen2.5-3B-Instruct 作为基座模型，该模型在指令理解与逻辑推理任务中表现出较好的稳定性与泛化能力。在参数高效微调框架下，系统采用低秩适配方法对模型进行微调，以降低训练成本并提升任务适配能力。

在具体配置方面，低秩适配的秩设置为 16，缩放系数设置为 16，随机失活率设置为 0。适配模块覆盖模型中的主要线性投影层，包括注意力与前馈网络中的关键映射层，从而增强模型对路由判别任务的表达能力。训练过程基于高效微调框架实现，并结合强化学习训练库进行优化，以提升训练速度并降低显存占用。

在优化策略方面，本研究采用自适应权重衰减优化算法，并通过低精度优化器进一步减少优化器状态的显存消耗。初始学习率设置为 $1\times10^{-4}$，并采用余弦退火策略进行动态调整，同时引入预热机制以提升训练稳定性。训练过程中单设备批次大小设置为 4，通过梯度累积将有效批次规模扩展至 32，并在较小规模数据集上限制训练轮次为 2，以降低过拟合风险。

\subsubsection{生成模型与检索后端}

在端到端检索增强生成框架中，生成模型与检索模块协同工作，共同构成系统的核心计算链路。生成模块选用 Meta-Llama-3-8B-Instruct 作为主要推理模型，并采用半精度模式进行加载。推理阶段的温度参数设置为 1.0，最大生成长度设置为 200，以兼顾生成结果的多样性与表达简洁性。

检索模块基于 Elasticsearch 构建，并融合了多种检索技术以提升召回能力与语义覆盖度。系统同时引入词法检索、语义检索与稀疏表示检索三种机制。词法检索采用 BM25 算法，用于捕捉查询中的关键词与实体信息；语义检索基于句向量模型构建稠密表示，并通过近似最近邻索引实现高效检索；稀疏表示检索通过学习稀疏向量形式，弥补词法检索与稠密表示在语义表达与精确匹配方面的不足。系统在检索阶段并行执行多路检索，并对不同检索结果进行加权融合，从而构建用于生成的上下文信息。

\subsubsection{硬件环境}

模型训练与推理评估均在高性能计算服务器上完成。服务器配备高性能多核处理器，以支撑数据预处理与检索服务的并发需求。深度学习相关任务在配备 32GB 显存的 GPU 环境中运行，并支持混合精度计算，从而满足大模型推理与中等规模模型微调的计算需求。该硬件配置为实验的稳定运行与结果复现提供了可靠保障。


\subsection{评估指标}
\label{sec:evaluation_metrics}

为了全面验证本文所提出的动态路由检索增强生成系统的有效性，本研究从分类器决策性能、端到端问答质量以及系统运行效率三个层面构建了多维度评估指标体系，以刻画系统在准确性、鲁棒性与计算开销之间的综合表现。

\subsubsection{分类性能指标}

路由分类器的判别能力直接影响下游检索策略的选择合理性，从而对整体系统性能产生关键影响。本研究围绕零检索、单跳检索与多跳推理三类决策结果对分类性能进行评估。设测试集样本总数为 $N$，类别集合为 $C=\{Z,S,M\}$，第 $i$ 个样本的真实标签为 $y_i$，预测标签为 $\hat{y}_i$。

首先，通过准确率指标衡量分类器在整体样本上的预测正确比例，其定义为
\begin{equation}
\text{Accuracy} = \frac{1}{N} \sum_{i=1}^{N} \mathbb{I}(y_i = \hat{y}_i),
\end{equation}
其中 $\mathbb{I}(\cdot)$ 为指示函数，当条件成立时取值为 1，否则取值为 0。该指标反映了分类器在整体层面的判别能力。

考虑到不同复杂度问题在实际场景中的分布可能存在显著差异，仅依赖准确率难以全面刻画模型性能，因此进一步引入基于类别的 F1 指标。对于任意类别 $c \in C$，定义其精确率与召回率为
\begin{equation}
P_c = \frac{TP_c}{TP_c + FP_c}, \quad R_c = \frac{TP_c}{TP_c + FN_c},
\end{equation}
其中 $TP_c$、$FP_c$ 与 $FN_c$ 分别表示类别 $c$ 的真阳性、假阳性与假阴性样本数。对应的 F1 值定义为
\begin{equation}
F1_c = \frac{2 \cdot P_c \cdot R_c}{P_c + R_c}.
\end{equation}
在此基础上，采用宏平均方式对各类别 F1 值进行汇总，得到综合指标
\begin{equation}
\text{Macro-F1} = \frac{1}{|C|} \sum_{c \in C} F1_c,
\end{equation}
以衡量分类器在不同类别上的整体平衡性。

% todo 三分类？

此外，为了分析分类错误的具体分布模式，本研究构建了三分类的混淆矩阵 $\mathbf{M}$，其中元素 $M_{jk}$ 表示真实类别为 $j$ 且被预测为类别 $k$ 的样本数量，从而揭示不同类别之间的混淆关系。

\subsubsection{问答性能指标}

针对生成结果的质量评估，本研究采用机器阅读理解任务中常用的指标体系。设标准化函数为 $\text{Normalize}(\cdot)$，用于对文本进行统一处理，包括大小写转换、标点符号与冗余空白的去除。对于第 $i$ 个样本，设预测答案为 $A_i$，标准答案集合为 $G_i=\{g_i^{(1)}, g_i^{(2)}, \ldots\}$。

首先，通过精确匹配指标衡量预测答案与标准答案在标准化后的完全一致性。该指标定义为
\begin{equation}
\text{EM}_i = \max_{g \in G_i} \mathbb{I}(\text{Normalize}(A_i) = \text{Normalize}(g)),
\end{equation}
最终结果为所有样本的平均值。该指标刻画了模型在严格匹配条件下的回答准确性。

其次，在词符层面引入 F1 指标以衡量预测答案与标准答案之间的语义重叠程度。将标准化后的预测答案与标准答案分别表示为词符集合 $T_A$ 与 $T_G$，定义共有词符集合为 $T_{common}=T_A \cap T_G$，则精确率与召回率分别为
\begin{equation}
P_{qa} = \frac{|T_{common}|}{|T_A|}, \quad R_{qa} = \frac{|T_{common}|}{|T_G|},
\end{equation}
对应的 F1 值定义为
\begin{equation}
\text{F1}_{qa} = \frac{2 \cdot P_{qa} \cdot R_{qa}}{P_{qa} + R_{qa}}.
\end{equation}
该指标能够更细粒度地刻画模型生成结果与标准答案之间的语义一致性。

考虑到生成式模型往往倾向于输出较为完整的自然语言表达，单纯依赖精确匹配指标可能导致对合理答案的低估，因此进一步引入基于包含关系的匹配指标。该指标定义为
\begin{equation}
\text{Acc}_i = \max_{g \in G_i} \left( \mathbb{I}(g \subseteq A_i) \lor \mathbb{I}(A_i \subseteq g) \right),
\end{equation}
其中 $\subseteq$ 表示字符串层面的包含关系。该指标用于缓解由于答案表达形式差异所导致的误判问题，从而更加合理地评估生成结果的有效性。

\subsubsection{效率指标}

动态路由机制的核心目标是在保证问答质量的同时降低系统计算开销，因此需要从检索开销与推理时间两个方面对系统效率进行量化评估。

首先，通过平均检索次数衡量系统在处理问题时对检索模块的调用频率。设第 $i$ 个样本的检索次数为 $k_i$，其中零检索策略对应 $k_i=0$，单跳检索策略对应 $k_i=1$，多跳推理策略对应 $k_i$ 位于预设区间内，则平均检索次数定义为
\begin{equation}
\frac{1}{N}\sum_{i=1}^{N} k_i.
\end{equation}
该指标反映了系统在不同路由决策下的检索负载水平。

其次，通过平均推理时间衡量系统从接收输入查询到生成最终答案的端到端耗时。该指标综合反映了分类、检索与生成三个阶段的时间成本，是评估系统整体效率的重要依据。

\section{实验结果和分析}\label{ch03: exp_result}

\subsection{路由分类器性能评估}
\label{sec:classifier_performance}

路由分类器作为动态检索增强生成系统中的核心决策模块，其主要功能是依据用户查询的语义特征与推理复杂度，预测所需的检索策略类型。分类器的判别能力直接影响系统在检索成本与回答准确性之间的权衡效果，即是否能够在避免冗余检索的同时，降低因检索不足而导致回答错误的风险。本节基于独立测试集，对路由分类器的整体性能与误判模式进行系统性评估。

\subsubsection{总体分类性能}

本研究在约 360 条样本构成的独立测试集上对分类器进行评估。测试集覆盖无检索、单跳检索与多跳推理三类问题复杂度，并保持与训练集相近的类别分布结构。

% TODO: 需要运行分类器评估实验，计算各类别(Z/S/M)的 Precision/Recall/F1
% 数据来源: evaluate/outputs/stage1_classifications/ 与真实标签对比
表 \ref{tab:classifier_metrics} 汇总了分类器在不同类别上的性能指标结果，包括分类准确性与类别级别的综合评价指标。实验结果表明，分类器在三分类任务上取得了较为稳定的性能表现，总体分类准确率达到 \textit{[请替换为真实准确率数值]}，宏平均指标达到 \textit{[请替换为真实数值]}。

\begin{table}[h]
\centering
\caption{路由分类器在测试集上的性能指标}
\label{tab:classifier_metrics}
% TODO: 需要运行分类器评估实验以获取以下数值
% 指标一=Precision, 指标二=Recall, 综合指标=F1
\begin{tabular}{lcccc}
\hline
类别 & 复杂度 & Precision（\%） & Recall（\%） & F1（\%） \\ \hline
Z 类别 & L0 & \textit{[请替换为真实数值]} & \textit{[请替换为真实数值]} & \textit{[请替换为真实数值]} \\
S 类别 & L1 & \textit{[请替换为真实数值]} & \textit{[请替换为真实数值]} & \textit{[请替换为真实数值]} \\
M 类别 & L2 & \textit{[请替换为真实数值]} & \textit{[请替换为真实数值]} & \textit{[请替换为真实数值]} \\ \hline
总体结果 & - & - & - & \textit{[请替换为真实数值]} \\
宏平均结果 & - & \textit{[请替换为真实数值]} & \textit{[请替换为真实数值]} & \textit{[请替换为真实数值]} \\ \hline
\end{tabular}
\end{table}

从表 \ref{tab:classifier_metrics} 可以观察到，分类器在无检索类别上的识别效果较为突出，其综合指标达到 \textit{[请替换为真实数值]}。这一现象表明，通过合成数据构造的训练策略能够有效覆盖逻辑推理、代码理解与闲聊等典型无检索场景，使模型能够较为清晰地区分内部知识与外部知识边界。在单跳检索与多跳推理两类复杂度较高的任务中，分类器同样表现出较强的区分能力，说明弱标签提示学习策略在复杂推理场景中具有较好的泛化效果。

\subsubsection{混淆矩阵与错误分析}

为了进一步刻画分类器的决策行为及其潜在风险，本研究构建了分类结果的混淆矩阵，如图 \ref{fig:confusion_matrix} 所示。通过分析不同类别之间的误判分布，可以揭示分类器在不同复杂度问题上的决策偏向，并评估其对下游检索增强生成系统的影响。

% TODO: 绘图 - 生成分类器三分类混淆矩阵热力图（Z/S/M三类）
% 依赖: 先运行分类器评估实验获取混淆矩阵数据
\begin{figure}[h]
\centering
\includegraphics[width=0.6\textwidth]{imgs/chapter03/confusion_matrix.pdf}
% TODO: 替换为实际生成的混淆矩阵热力图路径
\caption{路由分类器在测试集上的混淆矩阵}
\label{fig:confusion_matrix}
\end{figure}

% TODO: 以下比例数值依赖分类器评估实验的混淆矩阵结果
从混淆矩阵可以观察到，部分单跳检索问题被判定为多跳推理问题，其比例约为 \textit{[请替换为真实比例]}。在实际系统运行过程中，这意味着原本可以通过单次检索解决的问题被路由至多跳推理模块，从而引入额外的检索与推理开销。然而，由于多跳推理模块通常具备逐步验证与迭代检索机制，这类误判对最终回答准确性的影响相对有限，主要体现为效率层面的损耗。造成该现象的原因可能在于部分单跳问题包含复杂修饰语或低频实体，从而使模型倾向于采取更为谨慎的决策策略。

% TODO: 以下比例数值依赖分类器评估实验的混淆矩阵结果
相比之下，多跳推理问题被误判为单跳检索问题的比例相对较低，其占比约为 \textit{[请替换为真实比例]}。在检索增强生成系统中，这类误判可能导致系统无法获取完整的证据链，从而增加生成错误答案或出现幻觉的风险。实验结果表明，该类误判发生频率显著低于前述情形，说明分类器在训练过程中较好地捕捉了多跳推理问题的语义特征，例如跨文档关联、多实体关系与比较结构等，从而具备较强的风险规避能力。

涉及无检索类别的误判现象总体较少。少量无检索问题被判定为需要检索的问题，主要集中在常识与专业知识边界模糊的场景中，例如特定领域中的常识性问题。这类误判仅会引入有限的检索开销，对回答质量影响较小。另一方面，本研究几乎未观察到需要检索的问题被判定为无检索的问题，表明分类器能够有效避免闭卷回答所带来的严重幻觉风险。

综合上述分析可以发现，路由分类器不仅在整体性能指标上表现稳定，而且在错误分布上呈现出明显的非对称特征，即模型更倾向于采取保守策略，而非冒险性决策。这种错误分布特性与检索增强生成系统的工程设计原则高度一致，即在优先保障回答准确性的前提下，通过动态路由机制尽可能降低系统计算成本，为后续端到端实验提供了可靠的决策基础。

\subsection{端到端问答性能评估}\label{sec:e2e_results}

在完成分类器层面的性能分析后，本节进一步从端到端问答任务的角度，对所提出的动态路由方法在整体性能上的表现进行系统评估。实验将本文方法与三类基线系统进行对比，包括无检索生成（Z）、静态混合检索增强生成（S-Hybrid）以及复现的自适应检索增强生成方法（Adaptive-RAG）。所有方法均在六个基准数据集（约 3000 条测试样本）上进行评估，以确保实验结论的统计稳定性。

% TODO: 绘图 - 四种方法整体EM/F1对比柱状图

\subsubsection{整体性能对比}

表 \ref{tab:overall_comparison} 给出了四种方法在所有数据集上的宏观平均性能。实验结果表明，本文提出的动态路由方法在所有评估指标上均显著优于各基线系统。与无检索生成策略相比，本文方法在精确匹配率上提升了 18.46 个百分点，充分说明了外部知识检索对于事实型问答任务的重要性。与静态混合检索策略相比，本文方法的精确匹配率提升了 12.36 个百分点，表明自适应路由机制能够通过更合理的检索策略分配，显著降低过度检索与检索不足带来的性能损失。与复现的 Adaptive-RAG 方法相比，本文方法的精确匹配率仍高出 5.76 个百分点，验证了本文在训练数据构建与路由策略设计方面的改进效果。

\begin{table}[htbp]
    \centering
    \caption{四种方法在全部数据集上的整体性能对比（\%）}
    \label{tab:overall_comparison}
    \renewcommand{\arraystretch}{1.3}
    \begin{tabular}{l c c c c}
        \toprule
        \textbf{方法} & \textbf{EM} & \textbf{F1} & \textbf{Acc} & \textbf{Recall} \\
        \midrule
        无检索生成 (Z) & 1.00 & 2.20 & 1.90 & 2.63 \\
        静态混合检索 (S-Hybrid) & 7.10 & 15.33 & 24.23 & 26.50 \\
        Adaptive-RAG 基线 & 13.70 & 25.58 & 29.17 & 33.38 \\
        \textbf{本文方法} & \textbf{19.46} & \textbf{32.57} & \textbf{39.86} & \textbf{43.33} \\
        \bottomrule
    \end{tabular}
\end{table}

\subsubsection{分数据集性能分析}

为了进一步分析本文方法在不同复杂度问题上的表现差异，表 \ref{tab:proposed_per_dataset} 给出了本文方法在六个数据集上的详细性能分布。

\begin{table}[htbp]
    \centering
    \caption{本文方法在各数据集上的详细性能（\%）}
    \label{tab:proposed_per_dataset}
    \renewcommand{\arraystretch}{1.3}
    \begin{tabular}{l c c c c c}
        \toprule
        \textbf{数据集} & \textbf{EM} & \textbf{F1} & \textbf{Acc} & \textbf{Recall} & \textbf{样本数} \\
        \midrule
        NQ & 42.2 & 54.6 & 58.2 & 61.5 & 500 \\
        HotpotQA & 27.8 & 40.3 & 46.4 & 49.5 & 496 \\
        2WikiMultiHopQA & 14.4 & 28.4 & 44.4 & 46.7 & 500 \\
        TriviaQA & 13.6 & 36.2 & 40.0 & 47.2 & 500 \\
        SQuAD & 13.2 & 24.5 & 33.8 & 36.7 & 500 \\
        MuSiQue & 5.6 & 11.5 & 16.4 & 18.4 & 500 \\
        \bottomrule
    \end{tabular}
\end{table}

从表 \ref{tab:proposed_per_dataset} 可以观察到，系统在 NQ 数据集上取得了最高的精确匹配率（42.2\%），这与 NQ 数据集中大量问题属于单跳事实型查询的特征一致，路由分类器能够准确识别其复杂度并触发相应的单跳检索策略。在多跳数据集中，HotpotQA 的表现显著优于 MuSiQue，这主要是因为 MuSiQue 所要求的推理链条更长、组合复杂度更高，对多跳推理策略提出了更严苛的要求。

% TODO: 绘图 - 六数据集四种方法EM分组柱状图

表 \ref{tab:per_dataset_em_comparison} 进一步展示了四种方法在各数据集上的精确匹配率对比，以揭示路由策略在不同任务类型上的差异化表现。

\begin{table}[htbp]
    \centering
    \caption{四种方法在各数据集上的精确匹配率对比（EM, \%）}
    \label{tab:per_dataset_em_comparison}
    \renewcommand{\arraystretch}{1.3}
    \begin{tabular}{l c c c c}
        \toprule
        \textbf{数据集} & \textbf{Z} & \textbf{S-Hybrid} & \textbf{Adaptive-RAG} & \textbf{本文方法} \\
        \midrule
        SQuAD & 6.0 & 10.2 & 14.2 & 13.2 \\
        HotpotQA & 0.0 & 7.0 & 22.0 & 27.8 \\
        NQ & 0.0 & 15.6 & 20.8 & 42.2 \\
        TriviaQA & 0.0 & 7.2 & 10.0 & 13.6 \\
        MuSiQue & 0.0 & 1.0 & 2.4 & 5.6 \\
        2WikiMultiHopQA & 0.0 & 1.6 & 12.8 & 14.4 \\
        \bottomrule
    \end{tabular}
\end{table}

从表 \ref{tab:per_dataset_em_comparison} 的结果可以发现若干值得关注的现象。首先，在 NQ 数据集上，本文方法相对于 Adaptive-RAG 基线取得了最大幅度的提升（+21.4 个百分点），表明本文的增强型分类器在处理该数据集中的开放域事实查询时，能够更准确地选择最优检索策略。其次，在多跳数据集 HotpotQA 和 2WikiMultiHopQA 上，本文方法同样表现出明显优势，说明增强训练数据策略使分类器能够更准确地识别多跳推理需求，从而将这类问题正确路由至多跳推理模块。值得注意的是，在 SQuAD 数据集上，本文方法的精确匹配率略低于 Adaptive-RAG 基线（13.2\% vs 14.2\%），这可能与部分 SQuAD 样本被本文分类器路由至零检索路径有关，后续消融实验将对此进行进一步分析。

\subsubsection{消融实验：单策略性能分析}

为了分析路由机制的价值并理解不同执行策略的独立性能，本研究设计了消融实验，将所有测试样本统一分配至单一策略进行处理。表 \ref{tab:ablation_single_strategy} 给出了五种单策略配置的整体性能。

% TODO: 绘图 - 五种单策略消融对比柱状图

\begin{table}[htbp]
    \centering
    \caption{单策略消融实验：各策略在全部数据集上的整体性能（\%）}
    \label{tab:ablation_single_strategy}
    \renewcommand{\arraystretch}{1.3}
    \begin{tabular}{l c c c c}
        \toprule
        \textbf{策略} & \textbf{EM} & \textbf{F1} & \textbf{Acc} & \textbf{Recall} \\
        \midrule
        Z（零检索） & 1.00 & 2.20 & 1.90 & 2.63 \\
        S-Sparse（稀疏检索） & 7.13 & 15.37 & 24.20 & 26.47 \\
        S-Dense（稠密检索） & 6.13 & 14.12 & 24.00 & 26.15 \\
        S-Hybrid（混合检索） & 7.10 & 15.33 & 24.23 & 26.50 \\
        M（多跳推理） & 14.63 & 28.73 & 38.70 & 42.40 \\
        \bottomrule
    \end{tabular}
\end{table}

从表 \ref{tab:ablation_single_strategy} 中可以得出以下几点发现。第一，零检索策略在事实型问答任务上的表现极为有限（EM 仅为 1.00\%），这进一步验证了参数化知识在回答时效性和特定领域事实问题方面的局限性。第二，三种单跳检索策略的整体性能接近（EM 在 6.13\%--7.13\% 之间），其中稀疏检索与混合检索略优于纯稠密检索，说明在本文所用的异构数据集上，词法匹配信号对于命中答案所在文档具有重要作用。第三，M 策略以 14.63\% 的精确匹配率显著优于所有单跳策略，表明多跳推理机制在处理复杂问题时具备不可替代的优势。第四，本文提出的动态路由方法（EM = 19.46\%）超越了最优单策略 M（EM = 14.63\%）4.83 个百分点，有力地证明了自适应路由机制相比于"一刀切"式策略的优越性。

\subsubsection{路由动作分布分析}

为了深入理解分类器的决策行为，表 \ref{tab:action_distribution} 展示了本文方法在各数据集上的路由动作分布统计。

% TODO: 绘图 - 路由动作分布堆叠柱状图（按数据集）

\begin{table}[htbp]
    \centering
    \caption{本文方法在各数据集上的路由动作分布}
    \label{tab:action_distribution}
    \renewcommand{\arraystretch}{1.3}
    \begin{tabular}{l c c c c c}
        \toprule
        \textbf{数据集} & \textbf{Z} & \textbf{S-Sparse} & \textbf{S-Dense} & \textbf{S-Hybrid} & \textbf{M} \\
        \midrule
        SQuAD & 4 & 372 & 26 & 10 & 88 \\
        HotpotQA & 3 & 39 & 0 & 3 & 451 \\
        NQ & 3 & 479 & 5 & 1 & 12 \\
        TriviaQA & 4 & 379 & 5 & 4 & 108 \\
        MuSiQue & 0 & 22 & 0 & 3 & 475 \\
        2WikiMultiHopQA & 0 & 8 & 0 & 0 & 492 \\
        \bottomrule
    \end{tabular}
\end{table}

从表 \ref{tab:action_distribution} 的统计结果可以看出，分类器的路由决策呈现出明显的数据集特异性。在单跳事实型数据集（NQ、SQuAD、TriviaQA）上，分类器将绝大多数问题路由至单跳检索策略，其中稀疏检索占据主导地位。这与这类数据集中大量问题涉及精确实体匹配的特征高度一致。在多跳推理数据集（HotpotQA、MuSiQue、2WikiMultiHopQA）上，分类器将超过 90\% 的问题路由至 M 策略，表明分类器能够有效识别复杂推理需求。值得注意的是，SQuAD 数据集中有 88 个样本被路由至 M 策略，这可能是因为部分 SQuAD 问题的语言结构与多跳问题存在表层相似性。零检索策略被触发的频率极低，这与评估数据集均为知识密集型问答任务的特点一致。

%\section{实验设计}
%
%\subsection{数据集与评估指标}
%
%为了全面验证自适应路由框架在处理不同复杂度查询时的鲁棒性与泛化能力，本研究选取了五个具有代表性的开放域问答（Open-domain QA）基准数据集，涵盖了从简单事实检索到复杂逻辑推理的广泛难度谱系。
%
%数据集配置主要分为单跳（Single-hop）与多跳（Multi-hop）两大类：
%\begin{itemize}
%    \item \textbf{单跳推理基准}：采用了 \textit{Natural Questions (NQ)} 与 \textit{TriviaQA}。NQ 数据集源自真实的 Google 搜索引擎日志，包含大量非结构化的短查询，侧重于考察系统对真实用户意图的理解能力；TriviaQA 则包含来自问答网站的复杂事实查询，要求系统具备从富文本文档中提取精确证据的能力。此类数据集主要用于评估 $L_1$ 级路由策略的有效性。
%    \item \textbf{多跳推理基准}：选取了 \textit{HotpotQA}、\textit{2WikiMultihopQA} 与 \textit{MuSiQue}。HotpotQA 专注于桥接（Bridge）与比较（Comparison）类型的推理；2WikiMultihopQA 进一步引入了结构化的推理约束；MuSiQue 则以其高难度的连贯性推理（Connected Reasoning）著称，极易诱发检索漂移。这些数据集构成了测试 $L_2$ 级多跳 Agent 性能的核心试金石。
%\end{itemize}
%
%评估指标体系旨在量化系统在“准确率-效率”权衡中的表现，具体包含三个维度：
%\begin{enumerate}
%    \item \textbf{分类器性能}：评估路由模块的判别能力。我们定义了\textbf{粗粒度准确率}（Coarse-grained Accuracy），用于衡量模型在 $\{Z, S, M\}$ 三大类基础意图上的分类精度；以及\textbf{细粒度准确率}（Fine-grained Accuracy），用于考察模型在具体的检索形态（如 $S_{\text{sparse}}$ vs $S_{\text{dense}}$）上的预测一致性。
%    \item \textbf{端到端生成质量}：采用标准的\textbf{精确匹配率 (Exact Match, EM)} 与 \textbf{F1 分数}。EM 衡量生成的答案是否与标准答案字面完全一致，而 F1 分数则在 Token 级别评估答案的召回率与精确率，更能反映长文本答案的语义覆盖度。
%    \item \textbf{系统开销}：为了量化系统的运行效率，我们记录了平均推理延迟（Latency per Query）与平均 Token 消耗量。这组指标能够直观地反映路由策略在降低计算成本、避免无效检索方面的实际收益。
%\end{enumerate}
%
%\subsection{基线方法}
%
%为了验证本研究所提出的自适应路由框架在效率与准确率上的综合优势，我们将其实验结果与三类代表性的基线模型进行对比分析。这些基线涵盖了从静态检索到动态反射的不同技术范式：
%
%\begin{description}
%    \item[Standard RAG (Static)] 
%    该方法代表了目前主流的“一刀切”式检索增强生成架构。在此设置中，不存在路由决策模块，系统对测试集中的所有查询 $q$ 均无差别地执行固定参数的检索操作（通常配置为稠密向量检索 top-$k$ 文档），随后进行上下文生成。该基线主要用于量化由于缺乏复杂度感知而导致的过度检索（针对简单问题）或检索能力不足（针对复杂问题）对端到端性能的基础影响。
%
%    \item[Rule-based Router (Heuristic)] 
%    该方法采用基于人工规则的启发式路由策略。系统根据查询中的显式关键词（如“代码”、“翻译”、“对比”等）或句法特征，通过预定义的决策树将问题映射至特定的检索动作。例如，包含比较级词汇的查询被强制路由至多跳策略。该基线旨在评估传统规则方法在处理语义模糊或隐式推理需求时的局限性，并作为可学习分类器的下界参考。
%
%    \item[Self-RAG (Reflective)] 
%    Self-RAG \cite{selfrag} 代表了当前最先进的自适应 RAG 框架之一。不同于本文提出的前置式路由（Pre-emptive Routing），Self-RAG 采用生成时的反射机制，通过训练大模型在解码过程中动态生成特殊的反思词元（Reflection Tokens，如 \texttt{[Retrieve]} 或 \texttt{[No Retrieve]}）来按需调用检索模块。与 Self-RAG 的对比将重点揭示本文所主张的“感知-执行”解耦范式在推理延迟控制与全局策略规划上的差异优势。
%\end{description}
%
%\subsection{索引构建与检索流程}
%
%为了支撑自适应路由框架在不同粒度与模态下的检索需求，本实验构建了一个异构的多模态索引系统。该系统并未采用单一的向量检索方案，而是整合了稀疏倒排索引与稠密向量索引，形成了互补的检索视图。具体而言，索引架构包含三个核心组件：
%1. 稀疏索引 (Sparse Index)：采用经典的 BM25 算法捕捉基于词频的精确匹配信号，同时引入 SPLADE（Sparse Lexical and Expansion Model）生成学习型稀疏向量。SPLADE 通过在词表空间进行查询扩展，有效缓解了传统稀疏检索的“词汇不匹配”问题。
%2. 密集索引 (Dense Index)：利用预训练的 Transformer 编码器将文本映射为低维稠密向量，并基于 HNSW（Hierarchical Navigable Small World）算法构建近似最近邻图索引，以支持毫秒级的语义相似度检索。
%3. 重排序模块 (Reranking)：作为混合检索的收敛点，部署 Cross-Encoder 模型对召回的候选文档进行细粒度相关性打分。
%
%索引构建流程是一个计算密集型的工程过程。首先，对源语料库 $\mathcal{D}$ 执行清洗与分块（Chunking）预处理，确保文本片段长度与编码器的上下文窗口适配。随后，并行执行向量化推理与倒排表构建。为了平衡存储开销与检索效率，我们对稠密向量进行了量化处理，并针对 SPLADE 的稀疏性特征进行了存储优化。在实验过程中，索引的构建与维护占据了显著的计算资源，特别是在处理 WikiMultihop 等大规模语料时，需保证不同模态索引在文档 ID 层面的一致对齐。
%
%在检索执行阶段，系统根据分类器输出的 Action 动态选择最优的索引组合策略。对于被路由为 $S_{\text{sparse}}$ 的查询，系统仅访问倒排索引以最小化 I/O 开销；对于 $S_{\text{hybrid}}$ 及 $M$ 类动作，系统则并行触发稀疏与密集检索，并通过倒数排名融合（RRF）或加权线性融合聚合结果，随后进入重排序阶段。这种多维度的索引设计不仅为单一策略的性能评估提供了受控的实验环境，更为验证“不同复杂度问题适配不同检索模态”这一核心假设提供了坚实的物理基础。
%
%\subsection{路由分类器训练与优化}
%
%路由分类器 $\mathcal{M}_\theta$ 作为自适应框架的中枢决策单元，其训练过程不仅要求模型具备深度的语义理解能力，还需建立从自然语言表述到抽象检索策略的精准映射。为此，本研究在数据构建与模型微调层面实施了针对性的优化策略。
%
%在数据增强方面，鉴于现有的问答数据集缺乏显式的“检索策略”标签，我们提出了一种基于弱标签提示的知识蒸馏方案。具体而言，我们利用数据集的元信息作为弱监督信号，构建提示模版引导高参数量的教师模型（Teacher Model）生成结构化的思维链（Chain-of-Thought）分析与动作标签。该过程将隐含的数据集分布特征转化为显式的监督样本 $(q, a_{gold}, \text{reasoning})$，从而有效解决了训练数据的冷启动问题，并增强了模型对复杂逻辑的解释能力。
%
%在模型训练架构上，为了平衡训练效率与模型的通用语言能力，我们采用了**参数高效微调（Parameter-Efficient Fine-Tuning, PEFT）**技术，具体实现为低秩适应（Low-Rank Adaptation, LoRA）。我们保持基座大模型（Backbone LLM）的预训练权重 $W_0 \in \mathbb{R}^{d \times k}$ 冻结，仅在特定的注意力层注入可训练的低秩矩阵 $A \in \mathbb{R}^{d \times r}$ 和 $B \in \mathbb{R}^{r \times k}$，其中秩 $r \ll \min(d, k)$。前向传播更新公式如下：
%\begin{equation}
%h = W_0 x + \frac{\alpha}{r} BA x
%\end{equation}
%通过这种方式，模型能够以极小的参数增量（通常小于总参数量的 1\%）快速适配路由任务的特定分布，同时有效规避了全量微调中常见的灾难性遗忘（Catastrophic Forgetting）风险，确保模型在处理 $L_0$ 类通用任务时依然保持其基础的推理与生成能力。
%
%在工程优化层面，为了在大规模合成数据上实现稳定收敛，我们引入了混合精度训练（Mixed Precision Training）与梯度累积（Gradient Accumulation）策略。这不仅显著降低了显存占用，还通过模拟大批量数据（Large Batch Size）的梯度更新，平滑了损失函数的优化曲面，提升了分类器在极度不平衡类别分布下的泛化性能。
%
%\subsection{端到端实验设置}
%
%为了全面评估路由策略在实际问答场景中的效能，本实验构建了一个模块化的端到端推理流水线，根据分类器输出的决策动作 $a \in \mathcal{A}$ 动态调度异构的推理配置。具体执行逻辑如下：针对 $Z$ 类动作，系统旁路检索模块，直接调用大语言模型生成响应；针对 $S$ 类动作（$S_{\text{sparse}}, S_{\text{dense}}, S_{\text{hybrid}}$），系统首先并行执行对应的检索与重排序流程，构建增强上下文 $C$，随后进行一次性生成；针对 $M$ 类动作，系统实例化一个多跳推理 Agent，允许其在逻辑规划的指导下进行多轮次的“检索-阅读”迭代，直至生成最终答案。
%
%为了应对大规模测试集（如 HotpotQA 和 MuSiQue）带来的计算负载，并模拟高并发的生产环境，我们在推理阶段引入了异步并行与批量化处理机制。对于分类阶段与单跳生成阶段，采用向量化批量推理（Vectorized Batch Inference）以最大化 GPU 吞吐量；对于涉及复杂状态管理的多跳 Agent 推理，则采用多线程并发调度策略，显著降低了总体实验耗时。
%
%实验结果的量化评估主要基于生成的预测答案 $y_{\text{pred}}$ 与标准金标答案 $y_{\text{gold}}$ 之间的匹配程度。我们采用以下核心指标：
%\begin{itemize}
%    \item \textbf{精确匹配率 (Exact Match, EM)}：衡量 $y_{\text{pred}}$ 是否在规范化后与 $y_{\text{gold}}$ 实现字符串级的严格全等，反映了系统在事实性问题上的精确度。
%    \item \textbf{F1 分数 (F1 Score)}：计算预测答案与金标答案在 Token 级别的重叠程度，兼顾了召回率与精确率，适用于评估长文本生成的语义完整性。
%    \item \textbf{准确率 (Accuracy, ACC)}：用于衡量分类器本身做出正确路由决策的比例，作为系统性能归因的中间指标。
%\end{itemize}